\section*{Chapitre 3 \textemdash\ Calcul dans $\mathbb{R}$}
\addcontentsline{toc}{section}{Chapitre 3 \textemdash\ Calcul dans $\mathbb{R}$}

\subsection*{Exercices}
\addcontentsline{toc}{subsection}{Exercices}

%────────────────────────────────────────────────────────────────
\solutiontitle{3}{1}{Classification des nombres}

\begin{multicols}{3}
\begin{enumerate}
  \item $-7\in\Z$ (entier relatif, non naturel)
  \item $5/3\in\Q\setminus\Z$
  \item $\sqrt{16}=4\in\N$
  \item $\sqrt{3}\in\R\setminus\Q$
  \item $0\in\N$
  \item $\pi\in\R\setminus\Q$
  \item $-2/3\in\Q\setminus\Z$
  \item $\sqrt{9/4}=3/2\in\Q\setminus\Z$
  \item $\sqrt{2}\times\sqrt{8}=\sqrt{16}=4\in\N$
  \item $1{,}\overline{3}=4/3\in\Q\setminus\Z$
  \item $e\in\R\setminus\Q$
  \item $\sqrt{25}-5=0\in\N$
\end{enumerate}
\end{multicols}

%────────────────────────────────────────────────────────────────
\solutiontitle{3}{2}{Puissances~: calcul direct}

\begin{multicols}{3}
\begin{enumerate}
  \item $3^4\times3^{-2}=3^{4-2}=9$
  \item $2^7/2^3=2^4=16$
  \item $(5^2)^3=5^6=15625$
  \item $(-2)^5=-32$
  \item $4^{-2}=1/16$
  \item $\left(\tfrac{2}{3}\right)^3=8/27$
  \item $2^3\times5^3=(2\times5)^3=1000$
  \item $(3^2)^0=1$
  \item $\dfrac{6^4}{2^4}=\left(\dfrac{6}{2}\right)^4=81$
\end{enumerate}
\end{multicols}

%────────────────────────────────────────────────────────────────
\solutiontitle{3}{3}{Racines carr\'ees~: simplification}

\begin{multicols}{3}
\begin{enumerate}
  \item $\sqrt{72}=\sqrt{36\times2}=6\sqrt{2}$
  \item $\sqrt{50}/\sqrt{2}=\sqrt{25}=5$
  \item $(\sqrt{3}+1)(\sqrt{3}-1)=3-1=2$
  \item $\sqrt{12}+\sqrt{3}=2\sqrt{3}+\sqrt{3}=3\sqrt{3}$
  \item $(\sqrt{5})^4=25$
  \item $\sqrt{0{,}25}=0{,}5$
  \item $\sqrt{18}\times\sqrt{2}=\sqrt{36}=6$
  \item $\sqrt{75}-\sqrt{27}=5\sqrt{3}-3\sqrt{3}=2\sqrt{3}$
  \item $\dfrac{\sqrt{8}}{\sqrt{2}}=\sqrt{4}=2$
\end{enumerate}
\end{multicols}

%────────────────────────────────────────────────────────────────
\solutiontitle{3}{4}{Valeur absolue~: calcul}

\begin{multicols}{3}
\begin{enumerate}
  \item $\abs{-5}=5$
  \item $\abs{3-7}=4$
  \item $\abs{\pi-4}=4-\pi$ (car $\pi<4$)
  \item $\abs{-3}+\abs{2}=5$
  \item $\abs{1-2}=1$
  \item $\abs{5-2}=3$
  \item $\abs{2(-1)+1}=1$
  \item $\abs{(-3)^2-10}=1$
  \item $\bigl|\abs{5}-\abs{-3}\bigr|=2$
\end{enumerate}
\end{multicols}

%────────────────────────────────────────────────────────────────
\solutiontitle{3}{5}{Intervalles~: \'ecriture et repr\'esentation}

\begin{enumerate}
  \item Sur la droite~: $[-2,3[$ = crochet ferm\'e en $-2$, ouvert en $3$~; $]1,5]$~; $]-\infty,2[$~; $[0,+\infty[$.
  \item $[-1,4]$~; $]3,+\infty[$~; $]-\infty,-2]\cup[1,+\infty[$.
  \item $5\in[-1,5]$ est vraie~; $0\in]0,3[$ est fausse~;
        $-1\in]-2,0[$ est vraie.
\end{enumerate}

%────────────────────────────────────────────────────────────────
\solutiontitle{3}{6}{\'Ecriture scientifique}

\begin{multicols}{2}
\begin{enumerate}
  \item $45\,000\,000=4{,}5\times10^7$
  \item $0{,}000038=3{,}8\times10^{-5}$
  \item $3{,}14\times10^2\times5\times10^3=1{,}57\times10^6$
  \item $\dfrac{6\times10^8}{2\times10^{-3}}=3\times10^{11}$
  \item $0{,}0000001=1\times10^{-7}$
  \item $123\,456=1{,}23456\times10^5$
  \item $5{,}2\times10^3+3\times10^2=5{,}5\times10^3$
  \item $(2\times10^4)^3=8\times10^{12}$
\end{enumerate}
\end{multicols}

%────────────────────────────────────────────────────────────────
\solutiontitle{3}{7}{\'Equations simples avec valeur absolue}

\begin{enumerate}
  \item $\abs{x}=3 \Leftrightarrow x=3$ ou $x=-3$. Solution~: $\{-3,3\}$.
  \item $\abs{x}<5 \Leftrightarrow -5<x<5$. Solution~: $]-5,5[$.
  \item $\abs{x-2}=4 \Leftrightarrow x-2=4$ ou $x-2=-4$, soit $x=6$ ou $x=-2$.
  \item $\abs{x+1}\leq3 \Leftrightarrow -3\leq x+1\leq3$,
        soit $-4\leq x\leq2$. Solution~: $[-4,2]$.
  \item $\abs{x}>2$. Solution~: $]-\infty,-2[\cup]2,+\infty[$.
  \item $\abs{3x}=6\Leftrightarrow\abs{x}=2$. Solution~: $\{-2,2\}$.
  \item $\abs{x-1}=0\Leftrightarrow x=1$.
  \item $\abs{x+2}\geq0$ pour tout $x\in\R$. Solution~: $\R$.
\end{enumerate}

%────────────────────────────────────────────────────────────────
\solutiontitle{3}{8}{Op\'erations sur les fractions}

\begin{multicols}{2}
\begin{enumerate}
  \item $\dfrac{2}{3}+\dfrac{5}{6}=\dfrac{4}{6}+\dfrac{5}{6}=\dfrac{9}{6}=\dfrac{3}{2}$
  \item $\dfrac{3}{4}\times\dfrac{8}{9}=\dfrac{24}{36}=\dfrac{2}{3}$
  \item $\dfrac{5}{6}-\dfrac{1}{4}=\dfrac{10}{12}-\dfrac{3}{12}=\dfrac{7}{12}$
  \item $\dfrac{7}{3}\div\dfrac{14}{9}=\dfrac{7}{3}\times\dfrac{9}{14}=\dfrac{3}{2}$
  \item $\dfrac{1}{2}+\dfrac{1}{3}+\dfrac{1}{6}=1$
  \item $\left(\dfrac{3}{5}\right)^2-\dfrac{1}{25}=\dfrac{9}{25}-\dfrac{1}{25}=\dfrac{8}{25}$
\end{enumerate}
\end{multicols}

%────────────────────────────────────────────────────────────────
\solutiontitle{3}{9}{Intersection et r\'eunion d'intervalles}

\begin{enumerate}
  \item $A\cap B=[-3,5]\cap[1,8]=[1,5]$~; $A\cup B=[-3,8]$
  \item $A\cap B=]-2,4[\cap[3,7[=[3,4[$~; $A\cup B=]-2,7[$
  \item $A\cap B=]-\infty,2[\cap[1,+\infty[=[1,2[$~; $A\cup B=\R$
  \item $A\cap B=[-1,1]\cap[3,5]=\varnothing$~; $A\cup B=[-1,1]\cup[3,5]$
  \item $A\cap B=[0,4]\cap\R=[0,4]$~; $A\cup B=\R$
\end{enumerate}

%────────────────────────────────────────────────────────────────
\solutiontitle{3}{10}{Signe d'une expression}

\begin{enumerate}
  \item $x-3$~: positif si $x>3$, nul si $x=3$, n\'egatif si $x<3$.
  \item $2-x$~: positif si $x<2$, nul si $x=2$, n\'egatif si $x>2$.
  \item $x+5$~: positif si $x>-5$, nul si $x=-5$, n\'egatif si $x<-5$.
  \item $x^2$ est positif si $x\ne0$ et nul si $x=0$.
  \item $-x^2-1<0$ pour tout $x\in\R$.
\end{enumerate}

%────────────────────────────────────────────────────────────────
\solutiontitle{3}{11}{Encadrements}

\begin{multicols}{2}
\begin{enumerate}
  \item $4<\sqrt{17}<5$ car $4^2=16<17<25=5^2$
  \item $7<\sqrt{50}<8$ car $49<50<64$
  \item $2<\sqrt{2}+1<3$ car $1<\sqrt{2}<2$
  \item $4<2\sqrt{5}<5$ car $16<20<25$
  \item $9<\sqrt{99}<10$ car $81<99<100$
  \item $0<\sqrt{3}-1<1$ car $1<\sqrt{3}<2$
\end{enumerate}
\end{multicols}

%────────────────────────────────────────────────────────────────
\solutiontitle{3}{12}{Comparaison des r\'eels}

\begin{enumerate}
  \item $\sqrt{2}\approx1{,}414$, $3/2=1{,}5$, $1{,}4$, $\sqrt{1{,}9}\approx1{,}378$.
        Ordre~: $1{,}4 < \sqrt{1{,}9} < \sqrt{2} < 3/2$.
  \item $2\sqrt{3}\approx3{,}46$, $3\sqrt{2}\approx4{,}24$, $4$.
        Comparer~: $(2\sqrt{3})^2=12$, $(4)^2=16$, $(3\sqrt{2})^2=18$.
        Ordre~: $2\sqrt{3}<4<3\sqrt{2}$.
  \item $5/3\approx1{,}67$ vs $\sqrt{3}\approx1{,}73$~: $5/3<\sqrt{3}$
        (v\'erifier~: $(5/3)^2=25/9\approx2{,}78<3$).
\end{enumerate}

%────────────────────────────────────────────────────────────────
\solutiontitle{3}{13}{Racines~: d\'evelopper et factoriser}

\begin{multicols}{2}
\begin{enumerate}
  \item $(\sqrt{5}+\sqrt{2})(\sqrt{5}-\sqrt{2})=5-2=3$
  \item $(\sqrt{3}+1)^2=3+2\sqrt{3}+1=4+2\sqrt{3}$
  \item $(2-\sqrt{7})(2+\sqrt{7})=4-7=-3$
  \item $(\sqrt{2}+\sqrt{3})^2=2+2\sqrt{6}+3=5+2\sqrt{6}$
  \item $\sqrt{2}(\sqrt{8}-\sqrt{2})=\sqrt{16}-2=4-2=2$
  \item $(\sqrt{5}-\sqrt{3})^2=5-2\sqrt{15}+3=8-2\sqrt{15}$
\end{enumerate}
\end{multicols}

%────────────────────────────────────────────────────────────────
\solutiontitle{3}{14}{Lire un intervalle sur une droite}

\begin{enumerate}
  \item $2x-4>0\Leftrightarrow x>2$. Solution~: $]2,+\infty[$.
  \item $3x+6\leq0\Leftrightarrow x\leq-2$. Solution~: $]-\infty,-2]$.
  \item $-x+1\geq0\Leftrightarrow x\leq1$. Solution~: $]-\infty,1]$.
  \item $|x-3|<2\Leftrightarrow1<x<5$. Solution~: $]1,5[$.
\end{enumerate}

%────────────────────────────────────────────────────────────────
\solutiontitle{3}{15}{Puissances n\'egatives et fractionnaires}

\begin{multicols}{2}
\begin{enumerate}
  \item $\dfrac{x^3\cdot x^{-1}}{x^2}=\dfrac{x^2}{x^2}=1$ ($x\ne0$)
  \item $\left(\dfrac{a^2}{b}\right)^{-3}=\dfrac{b^3}{a^6}$ ($a,b\ne0$)
  \item $x^{1/2}\cdot x^{3/2}=x^{1/2+3/2}=x^2$ ($x\geq0$)
  \item $\dfrac{(2x)^3}{4x^2}=\dfrac{8x^3}{4x^2}=2x$ ($x\ne0$)
  \item $\left(x^{1/3}\right)^6=x^2$
  \item $a^{-1/2}=\dfrac{1}{\sqrt{a}}$ ($a>0$)
\end{enumerate}
\end{multicols}

%────────────────────────────────────────────────────────────────
\solutiontitle{3}{16}{In\'equations avec valeur absolue}

\begin{enumerate}
  \item $\abs{2x-3}\leq5\Leftrightarrow-5\leq2x-3\leq5\Leftrightarrow
        -2\leq2x\leq8\Leftrightarrow-1\leq x\leq4$.
        Solution~: $[-1,4]$.
  \item $\abs{x+4}>2\Leftrightarrow x+4>2$ ou $x+4<-2$,
        soit $x>-2$ ou $x<-6$.
        Solution~: $]-\infty,-6[\cup]-2,+\infty[$.
  \item $1\leq\abs{x-1}\leq4\Leftrightarrow$
        ($1\leq x-1\leq4$ ou $-4\leq x-1\leq-1$),
        soit $2\leq x\leq5$ ou $-3\leq x\leq0$.
        Solution~: $[-3,0]\cup[2,5]$.
  \item $\abs{3x+2}<\abs{x-1}$~: \'elever au carr\'e~:
        $(3x+2)^2<(x-1)^2$
        $\Rightarrow9x^2+12x+4<x^2-2x+1$
        $\Rightarrow8x^2+14x+3<0$
        $\Rightarrow(2x+3)(4x+1)<0$
        $\Rightarrow x\in]-3/2,-1/4[$.
\end{enumerate}

%────────────────────────────────────────────────────────────────
\solutiontitle{3}{17}{Rationalisations}

\begin{multicols}{2}
\begin{enumerate}
  \item $\dfrac{1}{\sqrt{3}}=\dfrac{\sqrt{3}}{3}$
  \item $\dfrac{2}{\sqrt{5}-1}=\dfrac{2(\sqrt{5}+1)}{4}=\dfrac{\sqrt{5}+1}{2}$
  \item $\dfrac{3}{\sqrt{7}+2}=\dfrac{3(\sqrt{7}-2)}{7-4}=\sqrt{7}-2$
  \item $\dfrac{\sqrt{3}+1}{\sqrt{3}-1}=\dfrac{(\sqrt{3}+1)^2}{2}=2+\sqrt{3}$
  \item $\dfrac{1}{\sqrt{2}+\sqrt{3}}=\dfrac{\sqrt{3}-\sqrt{2}}{3-2}=\sqrt{3}-\sqrt{2}$
  \item $\dfrac{6}{2-\sqrt{2}}=\dfrac{6(2+\sqrt{2})}{4-2}=3(2+\sqrt{2})$
\end{enumerate}
\end{multicols}

%────────────────────────────────────────────────────────────────
\solutiontitle{3}{18}{Preuve d'in\'egalit\'e entre racines}

\begin{enumerate}
  \item $(\sqrt{2}+\sqrt{3})^2=5+2\sqrt{6}$ et $(\sqrt{5})^2=5$.
        Comme $\sqrt{6}>0$~: $5+2\sqrt{6}>5$, donc $\sqrt{2}+\sqrt{3}>\sqrt{5}$.
        En particulier, $\sqrt{2}+\sqrt{3}\ne\sqrt{5}$.\hfill$\square$
  \item Voir ci-dessus~: $\sqrt{2}+\sqrt{3}>\sqrt{5}$.\hfill$\square$
  \item $(\sqrt{n+1}-\sqrt{n})(\sqrt{n+1}+\sqrt{n})=1$,
        donc $\sqrt{n+1}-\sqrt{n}=\dfrac{1}{\sqrt{n+1}+\sqrt{n}}$.
        Comme $\sqrt{n+1}+\sqrt{n}>2\sqrt{n}>0$~:
        $0<\sqrt{n+1}-\sqrt{n}<\dfrac{1}{2\sqrt{n}}\to0$.\hfill$\square$
\end{enumerate}

%────────────────────────────────────────────────────────────────
\solutiontitle{3}{19}{Signe et tableau}

\begin{enumerate}
  \item $(x-2)(x+3)$~: racines $x=2$ et $x=-3$.
\begin{center}\renewcommand{\arraystretch}{1.3}
\begin{tabular}{|c|ccccccc|}
\hline $x$ & $-\infty$ & & $-3$ & & $2$ & & $+\infty$ \\\hline
$(x-2)(x+3)$ & & $+$ & $0$ & $-$ & $0$ & $+$ & \\\hline
\end{tabular}
\end{center}
  \item $(2x-1)(-x+4)$~: racines $x=1/2$ et $x=4$.
        Positif sur $[1/2,4]$.
  \item $\dfrac{x+1}{x-3}$~: nul en $x=-1$, interdit en $x=3$.
        Positif sur $]-\infty,-1]\cup]3,+\infty[$.
  \item $\dfrac{(x-1)(x+2)}{x-5}$~: nul en $1$ et $-2$, interdit en $5$.
        Positif sur $[-2,1]\cup]5,+\infty[$.
\end{enumerate}

%────────────────────────────────────────────────────────────────
\solutiontitle{3}{20}{\'Equations sans carr\'es}

\begin{enumerate}
  \item $\sqrt{x}=3 \Leftrightarrow x=9$ (et $9\geq0$~: valide). Solution~: $\{9\}$.
  \item $\sqrt{2x+1}=3 \Leftrightarrow 2x+1=9 \Leftrightarrow x=4$.
        V\'erif.~: $\sqrt{9}=3$.\checkmark\ Solution~: $\{4\}$.
  \item $\sqrt{x+5}=x-1$. Condition~: $x\geq1$ et $x+5\geq0$.
        $x+5=(x-1)^2=x^2-2x+1 \Rightarrow x^2-3x-4=0 \Rightarrow (x-4)(x+1)=0$.
        $x=4$ ou $x=-1$. V\'erif.~: $x=4$~: $3=3$\checkmark~; $x=-1$~: $\sqrt{4}=2\ne-2$, rejet\'e.
        Solution~: $\{4\}$.
  \item $\sqrt{x}+\sqrt{x+1}=\sqrt{4x-1}$.
        Domaine~: $x\geq0$ et $4x-1\geq0$, soit $x\geq1/4$.
        En \'elevant au carr\'e~: $x+(x+1)+2\sqrt{x(x+1)}=4x-1$
        $\Rightarrow 2\sqrt{x(x+1)}=2x-2 \Rightarrow \sqrt{x^2+x}=x-1$.
        Condition~: $x\geq1$. \'El.~: $x^2+x=(x-1)^2=x^2-2x+1 \Rightarrow 3x=1 \Rightarrow x=1/3$.
        Mais $1/3<1$~: rejet\'e. Pas de solution.
\end{enumerate}

%────────────────────────────────────────────────────────────────
\solutiontitle{3}{21}{Cha\^ines d'in\'egalit\'es}

\begin{enumerate}
  \item $a\leq b$ et $b\leq c \Rightarrow a\leq c$ (transitivit\'e de $\leq$).
        G\'en\'eralisation~: si $a_1\leq a_2\leq\cdots\leq a_n$, alors $a_1\leq a_n$.
  \item $a<b$ et $c>0 \Rightarrow ac<bc$ (multiplier par positif conserve le sens).
  \item $0<a<b \Rightarrow a^2<b^2$, car $a^2<ab<b^2$. R\'eciproque~: vraie pour $a,b>0$.
        Pour $a,b\in\R$~: $a^2<b^2\Leftrightarrow |a|<|b|$, pas forc\'ement $a<b$.
  \item Si $0<a\leq1$~: $a^n\leq a$ pour $n\geq1$ (car multiplier par $a\leq1$
        diminue). Si $a>1$~: $a^n>a$.
\end{enumerate}

%────────────────────────────────────────────────────────────────
\solutiontitle{3}{22}{Approximation de r\'eels}

\begin{enumerate}
  \item $2{,}23^2=4{,}9729<5$ et $2{,}24^2=5{,}0176>5$~:
        $2{,}23<\sqrt{5}<2{,}24$.
  \item $3{,}14<\pi<3{,}15$ (valeur connue).
  \item $\sqrt{2}\approx1{,}41421$. Fractions proches~: $99/70\approx1{,}41429$,
        $7/5=1{,}4$, $17/12\approx1{,}41\overline{6}$.
        La fraction $99/70$ approche $\sqrt{2}$ \`a $10^{-5}$ pr\`es.
\end{enumerate}

%────────────────────────────────────────────────────────────────
\solutiontitle{3}{23}{Distances sur la droite r\'eelle}

\begin{enumerate}
  \item $d(3,-5)=|3-(-5)|=8$.
  \item $\{x\mid|x-1|\leq2\}=[-1,3]$.
  \item $|x-a|<r\Leftrightarrow a-r<x<a+r$~: c'est l'intervalle $]a-r,a+r[$.
  \item Les points \'equidistants de $a$ et $b$ v\'erifient $|x-a|=|x-b|$,
        soit $x=(a+b)/2$~: le milieu du segment $[a,b]$.
\end{enumerate}

%────────────────────────────────────────────────────────────────
\solutiontitle{3}{24}{Fractions alg\'ebriques}

\begin{multicols}{2}
\begin{enumerate}
  \item $\dfrac{x^2-4}{x-2}=\dfrac{(x-2)(x+2)}{x-2}=x+2$ ($x\ne2$)
  \item $\dfrac{x^2-1}{x^2+2x+1}=\dfrac{(x-1)(x+1)}{(x+1)^2}=\dfrac{x-1}{x+1}$ ($x\ne-1$)
  \item $\dfrac{x^2+3x+2}{x+2}=\dfrac{(x+1)(x+2)}{x+2}=x+1$ ($x\ne-2$)
  \item $\dfrac{x^3-1}{x-1}=\dfrac{(x-1)(x^2+x+1)}{x-1}=x^2+x+1$ ($x\ne1$)
\end{enumerate}
\end{multicols}

%────────────────────────────────────────────────────────────────
\solutiontitle{3}{25}{Puissances et ordre de grandeur}

\begin{enumerate}
  \item $0{,}0001\text{ m}=10^{-4}\text{ m}$.
        Dans $1\text{ m}$~: $10^4=10\,000$ grains. Dans $1\text{ km}=10^3\text{ m}$~:
        $10^7$ grains.
  \item La Terre a un rayon $\approx6{,}4\times10^6\text{ m}$.
        Diam\`etre $\approx1{,}28\times10^7\text{ m}=1{,}28\times10^{10}\text{ mm}$.
        Grains en ligne~: $1{,}28\times10^{10}/0{,}1=1{,}28\times10^{11}$.
  \item $N_A\approx6{,}02\times10^{23}$~; $1\text{ mol}=6{,}02\times10^{23}$ mol\'ecules.
        Masse d'une mol\'ecule d'eau ($M=18\text{ g/mol}$)~:
        $18/(6{,}02\times10^{23})\approx3\times10^{-23}\text{ g}$.
\end{enumerate}

%────────────────────────────────────────────────────────────────
\solutiontitle{3}{26}{Valeur absolue et distance}

\begin{enumerate}
  \item $a\geq0$~: $|a|=a=|-a|$.\quad $a<0$~: $|a|=-a$ et $|-a|=|{-a}|=-(-a)=a=|a|$.
        Dans les deux cas $|a|=|-a|$.\hfill$\square$
  \item In\'egalit\'e triangulaire~: $|a+b|\leq|a|+|b|$.
        Cas~1~: $a+b\geq0$~: $|a+b|=a+b\leq|a|+|b|$ car $a\leq|a|$ et $b\leq|b|$.
        Cas~2~: $a+b<0$~: $|a+b|=-(a+b)=(-a)+(-b)\leq|a|+|b|$.\hfill$\square$
  \item $|a-b|$ est la distance entre $a$ et $b$ sur la droite.
        $|a-b|\leq|a|+|b|$ car $|a-b|=|a+(-b)|\leq|a|+|-b|=|a|+|b|$.
  \item \'Egalit\'e $|a+b|=|a|+|b|$ ssi $a$ et $b$ ont le m\^eme signe ou l'un est nul.
\end{enumerate}

%────────────────────────────────────────────────────────────────
\solutiontitle{3}{27}{Syst\`eme d'in\'equations}

\begin{enumerate}
  \item $\begin{cases}x+1>0\\2x-3\leq5\end{cases}\Leftrightarrow\begin{cases}x>-1\\x\leq4\end{cases}$
        Solution~: $]-1,4]$.
  \item $\begin{cases}3x-1>x+5\\x-2\leq3x+4\end{cases}\Leftrightarrow\begin{cases}x>3\\x\geq-3\end{cases}$
        Solution~: $]3,+\infty[$.
  \item $\begin{cases}x^2\geq1\\\abs{x}\leq3\end{cases}$~:
        $x^2\geq1\Leftrightarrow x\leq-1$ ou $x\geq1$~;
        $\abs{x}\leq3\Leftrightarrow-3\leq x\leq3$.
        Solution~: $[-3,-1]\cup[1,3]$.
\end{enumerate}

%────────────────────────────────────────────────────────────────
\solutiontitle{3}{28}{Encadrements et d\'ecimales}

\begin{enumerate}
  \item $1^3=1<2$ et $2^3=8>2$~: $1<\sqrt[3]{2}<2$.\hfill$\square$
  \item $1{,}25^3=1{,}953125<2$ et $1{,}26^3=2{,}000376>2$~:
        $1{,}25<\sqrt[3]{2}<1{,}26$.\hfill$\square$
  \item Dichotomie~: $m=1{,}255$, $1{,}255^3\approx1{,}977<2$,
        donc $1{,}255<\sqrt[3]{2}<1{,}260$.
        $m=1{,}259$~: $1{,}259^3\approx1{,}997<2$~; $m=1{,}2599$~: environ $2$.
        $\sqrt[3]{2}\approx1{,}2599$.
  \item En g\'en\'eral, la dichotomie donne $k$ d\'ecimales en $k$ it\'erations
        (la pr\'ecision double environ tous les 3{,}32 it\'erations).
\end{enumerate}

%────────────────────────────────────────────────────────────────
\solutiontitle{3}{29}{Algorithme~: encadrement de $\sqrt{n}$}

Trace pour $n=7$~:
$k=1$~: $1^2=1<7$~; $k=2$~: $4<7$~; $k=3$~: $9>7$~: \textbf{arr\^et}.
R\'esultat~: $2<\sqrt{7}<3$ (encadrement entier).

V\'erification~: $2{,}6^2=6{,}76<7<7{,}29=2{,}7^2$, donc $2{,}6<\sqrt{7}<2{,}7$.

L'invariant de boucle est~: \og apr\`es l'it\'eration $k$, on sait que $k^2\leq n$\fg.

%────────────────────────────────────────────────────────────────
\solutiontitle{3}{30}{In\'egalit\'e par carr\'es}

\begin{enumerate}
  \item $(a-b)^2\geq0 \Rightarrow a^2-2ab+b^2\geq0 \Rightarrow a^2+b^2\geq2ab$.\hfill$\square$
  \item Pour $a,b>0$~: $a^2+b^2\geq2ab$ et $ab>0$,
        donc $\dfrac{a^2+b^2}{2ab}\geq1$,
        d'o\`u $\dfrac{a}{b}+\dfrac{b}{a}\geq2$
        (in\'egalit\'e AM--GM appliqu\'ee \`a $a/b$ et $b/a$).
  \item \'Egalit\'e ssi $a=b$.\hfill$\square$
\end{enumerate}

%────────────────────────────────────────────────────────────────
\solutiontitle{3}{31}{Irrationnalit\'e de $\sqrt{3}$}

Supposons $\sqrt{3}=p/q$ avec $\pgcd(p,q)=1$.
Alors $p^2=3q^2$, donc $3\mid p^2$.
Si $p=3k$~: $9k^2=3q^2$, d'o\`u $3\mid q^2$, puis $3\mid q$.
Mais $3\mid\pgcd(p,q)$~: contradiction avec $\pgcd(p,q)=1$.\hfill$\square$

%────────────────────────────────────────────────────────────────
\solutiontitle{3}{32}{Puissances de $(1+x)$ et encadrements}

\begin{enumerate}
  \item $(1+x)^2=1+2x+x^2$~; $(1+x)^3=1+3x+3x^2+x^3$.
  \item $(1+x)^2=1+2x+x^2\geq1+2x$ car $x^2\geq0$.\hfill$\square$
  \item Plus g\'en\'eralement, pour $x\geq-1$ et $n\geq1$~:
        $(1+x)^n\geq1+nx$ (in\'egalit\'e de Bernoulli).
        Pour $n=3$~: $(1+x)^3=1+3x+3x^2+x^3$.
        Si $x\geq0$~: tous les termes sont positifs, $(1+x)^3\geq1+3x$.\checkmark
        Si $-1\leq x<0$~: \'etude plus d\'etaill\'ee.
\end{enumerate}

%────────────────────────────────────────────────────────────────
\solutiontitle{3}{33}{Somme de deux irrationnels}

\begin{enumerate}
  \item Non~: $\sqrt{2}+(-\sqrt{2})=0\in\Q$. Contre-exemple donn\'e.
  \item Non~: $\sqrt{2}\times\sqrt{2}=2\in\Q$.
  \item $(\sqrt{2}+1)(\sqrt{2}-1)=2-1=1\in\Q$~: produit de deux irrationnels rationnel.
  \item La diff\'erence d'un rationnel et d'un irrationnel est irrationnelle~:
        si $r\in\Q$ et $x\notin\Q$, alors $r-x\notin\Q$ (car si $r-x=s\in\Q$,
        alors $x=r-s\in\Q$, contradiction).\hfill$\square$
\end{enumerate}

%────────────────────────────────────────────────────────────────
\solutiontitle{3}{34}{In\'egalit\'e des moyennes}

\begin{enumerate}
  \item AM $\geq$ GM~: $\dfrac{a+b}{2}\geq\sqrt{ab}$ pour $a,b\geq0$.
        Preuve~: $(\sqrt{a}-\sqrt{b})^2\geq0
        \Rightarrow a-2\sqrt{ab}+b\geq0
        \Rightarrow\dfrac{a+b}{2}\geq\sqrt{ab}$.\hfill$\square$
  \item \'Egalit\'e ssi $\sqrt{a}=\sqrt{b}$, i.e., $a=b$.
  \item Application~: pour $a,b>0$ avec $a+b=S$ fix\'e~:
        $ab\leq\left(\dfrac{a+b}{2}\right)^2=\dfrac{S^2}{4}$.
        Le produit est maximum quand $a=b=S/2$.
  \item Exemple~: $a=3$, $b=5$~: AM~$=4$, GM~$=\sqrt{15}\approx3{,}87<4$.\checkmark
\end{enumerate}

%────────────────────────────────────────────────────────────────
\solutiontitle{3}{35}{Approximer un irrationnel par des rationnels}

\renewcommand{\arraystretch}{1.4}
\begin{center}
\begin{tabular}{|c|c|c|c|}
\hline
\rowcolor{BN!10}
$p/q$ & valeur & $|p/q-\sqrt{2}|$ & $q^2|p/q-\sqrt{2}|$ \\
\hline
$1/1$ & $1{,}000$ & $0{,}414$ & $0{,}414$ \\
$3/2$ & $1{,}500$ & $0{,}086$ & $0{,}344$ \\
$7/5$ & $1{,}400$ & $0{,}014$ & $0{,}354$ \\
$17/12$ & $1{,}4167$ & $0{,}003$ & $0{,}420$ \\
$41/29$ & $1{,}4138$ & $0{,}0004$ & $0{,}334$ \\
\hline
\end{tabular}
\end{center}

Chaque fois, la s\'equence $(p_n/q_n)$ v\'erifie la r\'ecurrence $p_{n+1}=p_n+2q_n$,
$q_{n+1}=p_n+q_n$ (algorithme des convergents).

%────────────────────────────────────────────────────────────────
\solutiontitle{3}{36}{Propri\'et\'es de la valeur absolue}

\begin{enumerate}
  \item $|ab|=|a|\cdot|b|$~:
        Si $a,b\geq0$~: $|ab|=ab=|a||b|$.\quad
        Si $a\geq0,b<0$~: $|ab|=|a\cdot b|=-ab=a(-b)=|a||b|$.
        (Les cas $a<0,b\geq0$ et $a,b<0$ sont sym\'etriques.)\hfill$\square$
  \item $\left|\dfrac{a}{b}\right|=\dfrac{|a|}{|b|}$ ($b\ne0$)~:
        $\left|\dfrac{a}{b}\right|=|a\cdot b^{-1}|=|a|\cdot|b^{-1}|=|a|/|b|$.\hfill$\square$
  \item $|a|^2=a^2$~: si $a\geq0$~: $|a|^2=a^2$.\quad si $a<0$~: $|a|^2=(-a)^2=a^2$.\hfill$\square$
  \item $||a|-|b||\leq|a-b|$~: in\'egalit\'e triangulaire appliqu\'ee \`a
        $|a|=|(a-b)+b|\leq|a-b|+|b|$, d'o\`u $|a|-|b|\leq|a-b|$.
        De m\^eme $|b|-|a|\leq|b-a|=|a-b|$. Donc $||a|-|b||\leq|a-b|$.\hfill$\square$
\end{enumerate}

%────────────────────────────────────────────────────────────────
\solutiontitle{3}{37}{Approximation de $\pi$}

Archim\`ede obtint $3\frac{10}{71}<\pi<3\frac{1}{7}$, soit
$3{,}1408\ldots<\pi<3{,}1428\ldots$

\begin{enumerate}
  \item Un polygone r\'egulier \`a $n$ c\^ot\'es inscrit dans le cercle unit\'e
        a pour p\'erim\`etre $P_n=2n\sin(\pi/n)$.
        Pour $n=96$~: $P_{96}=192\sin(\pi/96)\approx3{,}1410$.
  \item Le polygone circonscrit a pour p\'erim\`etre $P_n'=2n\tan(\pi/n)$.
        Pour $n=96$~: $P_{96}'\approx3{,}1427$.
  \item Archim\`ede a obtenu ces valeurs par calcul g\'eom\'etrique, sans
        calculatrice, en partant de $n=6$ et en doublant \`a chaque \'etape.
  \item $P_n\to2\pi$ car $n\sin(\pi/n)=\pi\cdot\dfrac{\sin(\pi/n)}{\pi/n}\to\pi$.
\end{enumerate}

%────────────────────────────────────────────────────────────────
\solutiontitle{3}{38}{\logiq\ Propositions sur les ensembles}

\begin{enumerate}
  \item $\forall x\in\N, x\in\Z$~: \textbf{Vraie} ($\N\subset\Z$).
        N\'egation~: $\exists x\in\N, x\notin\Z$. Fausse.
  \item $\exists x\in\Q, x^2=2$~: \textbf{Fausse} ($\sqrt{2}\notin\Q$).
        N\'egation~: $\forall x\in\Q, x^2\ne2$. Vraie.
  \item $\forall x\in\R, x^2\geq0$~: \textbf{Vraie}.
        N\'egation~: $\exists x\in\R, x^2<0$. Fausse.
  \item $\exists x\in\Z, x/2\in\Z$~: \textbf{Vraie} ($x=2$ convient).
        N\'egation~: $\forall x\in\Z, x/2\notin\Z$. Fausse.
\end{enumerate}

%────────────────────────────────────────────────────────────────
\solutiontitle{3}{39}{\logiq\ Implication et valeur absolue}

\begin{enumerate}
  \item Contrapos\'ee~: si $x\geq3$ ou $x\leq-3$, alors $|x|\geq3$.
        Preuve~: $|x|<3\Leftrightarrow-3<x<3$ (par d\'efinition). D'o\`u la contrapos\'ee.\hfill$\square$
  \item $|x-2|<1\Rightarrow1<x<3\Rightarrow x>0$~: vraie.\hfill$\square$
        Contrapos\'ee~: $x\leq0\Rightarrow|x-2|\geq1$.
        Si $x\leq0$~: $x-2\leq-2$, donc $|x-2|=2-x\geq2\geq1$.\hfill$\square$
  \item R\'eciproque de~(1)~: si $-3<x<3$, alors $|x|<3$. \textbf{Vraie}
        (c'est la d\'efinition~: $|x|<3\Leftrightarrow-3<x<3$, donc les deux sens).
\end{enumerate}

%────────────────────────────────────────────────────────────────
\solutiontitle{3}{40}{\logiq\ Nier une condition d'appartenance}

\begin{enumerate}
  \item $\forall x\in[-2,3], x^2\leq9$~: \textbf{Vraie}
        ($|x|\leq3$ sur $[-2,3]$ donc $x^2\leq9$).
        N\'egation~: $\exists x\in[-2,3], x^2>9$. Fausse.
  \item $\exists x\in\R, \sqrt{x}<0$~: \textbf{Fausse} ($\sqrt{x}\geq0$ toujours).
        N\'egation~: $\forall x\in\R, \sqrt{x}\geq0$. Vraie.
  \item $\forall x\in\R^*, 1/x>0$~: \textbf{Fausse} (contre-exemple~: $x=-1$).
        N\'egation~: $\exists x\in\R^*, 1/x\leq0$. Vraie.
  \item $\forall a,b\in\R, |a+b|=|a|+|b|$~: \textbf{Fausse} (contre-exemple~: $a=1, b=-1$).
        N\'egation~: $\exists a,b\in\R, |a+b|\ne|a|+|b|$. Vraie.
\end{enumerate}

%────────────────────────────────────────────────────────────────
\solutiontitle{3}{41}{\logiq\ Condition n\'ecessaire et suffisante}

\begin{enumerate}
  \item $x>0$ pour que $\sqrt{x}$ existe~: \textbf{condition suffisante} (pas CN~:
        $\sqrt{0}$ existe aussi, et $x=0$ convient). Correction~: $x\geq0$ est CNS.
  \item $x=0$ pour que $x^2=0$~: \textbf{CNS}. $x=0\Rightarrow x^2=0$ et $x^2=0\Rightarrow x=0$.\hfill$\square$
  \item $|x|<3$ pour que $x^2<9$~: \textbf{CNS}. $|x|<3\Leftrightarrow x^2<9$.\hfill$\square$
  \item $\Delta>0$ pour que $ax^2+bx+c=0$ ait deux racines~: \textbf{CNS} (d\'efinition).\hfill$\square$
\end{enumerate}

%────────────────────────────────────────────────────────────────
\solutiontitle{3}{42}{\logiq\ R\'ediger une d\'emonstration par l'absurde}

\begin{enumerate}
  \item Supposons qu'il existe un plus grand rationnel $M$.
        Alors $M+1$ est rationnel et $M+1>M$~: contradiction.\hfill$\square$
  \item Supposons $a^2$ pair et $a$ impair~: $a=2k+1$, $a^2=4k^2+4k+1$ impair~:
        contradiction.\hfill$\square$
  \item Supposons $r=\sqrt{2}+\sqrt{3}\in\Q$.
        Alors $r-\sqrt{2}=\sqrt{3}$ et $r^2=5+2\sqrt{6}$,
        d'o\`u $\sqrt{6}=(r^2-5)/2\in\Q$.
        Mais $\sqrt{6}$ est irrationnel (preuve analogue \`a $\sqrt{2}$)~: contradiction.\hfill$\square$
\end{enumerate}

%────────────────────────────────────────────────────────────────
\solutiontitle{3}{43}{\logiq\ Logique des in\'egalit\'es}

\begin{enumerate}
  \item \textbf{Fausse.} Contre-exemple~: $a=1>b=-2$, mais $a^2=1<4=b^2$.
  \item $a>b\Rightarrow a^2>b^2$ ssi $a>|b|$, i.e., $a>0$ et $a>-b$.
        Condition exacte~: $a>0$ et $b>-a$, i.e., $a>|b|$.
  \item $a^2>b^2\Rightarrow|a|>|b|$, pas forc\'ement $a>b$.
        Contre-exemple~: $a=-3, b=2$~: $9>4$ mais $-3<2$.
\begin{center}\renewcommand{\arraystretch}{1.3}
\begin{tabular}{|c|c|c|c|c|}
\hline
\rowcolor{BN!10}
$(a,b)$ & $a>b$~? & $a^2>b^2$~? & impl.~? \\ \hline
$(2,1)$ & Oui & Oui & \checkmark \\
$(-2,1)$ & Non & Oui & --- \\
$(2,-1)$ & Oui & Oui & \checkmark \\
$(-1,-2)$ & Oui & Non & $\times$ \\
\hline
\end{tabular}
\end{center}
\end{enumerate}

%────────────────────────────────────────────────────────────────
\solutiontitle{3}{44}{\algo\ Algorithme~: d\'ecimales de $\sqrt{n}$}

Apr\`es avoir trouv\'e $k$ tel que $k^2\leq n<(k+1)^2$~:
phase~2 avec pas $h=0{,}1$, phase~3 avec $h=0{,}01$, etc.

Pour $n=2$~: $k=1$. Phase~2~: $1{,}4^2=1{,}96<2<2{,}25=1{,}5^2$, d'o\`u
$1{,}4<\sqrt{2}<1{,}5$. Phase~3~: $1{,}41^2=1{,}9881<2<1{,}4\overline{4}^2$, d'o\`u
$1{,}41<\sqrt{2}<1{,}42$. Etc.

$\sqrt{2}\approx1{,}41421$.
Nombre d'it\'erations~: environ $10$ par d\'ecimale (10 + 10 + 10 + 10 + 10 = 50 totales).

%────────────────────────────────────────────────────────────────
\solutiontitle{3}{45}{\algo\ Algorithme~: test d'appartenance}

Tests~:
\begin{enumerate}
  \item $(a,b,x)=(1,5,3)$~: $3\geq1$ et $3\leq5$~: \Vrai.
  \item $(1,5,7)$~: $7\leq5$ est faux~: \Faux.
  \item $(1,5,1)$~: $1\geq1$ et $1\leq5$~: \Vrai.
\end{enumerate}

Versions modifi\'ees~:
$]a,b[$~: remplacer $\geq a$ et $\leq b$ par $>a$ et $<b$.
$[a,b[$~: $x\geq a$ \textbf{et} $x<b$.
$]a,b]$~: $x>a$ \textbf{et} $x\leq b$.

En Python~: \texttt{a <= x <= b}, \texttt{a < x < b}, etc.

%────────────────────────────────────────────────────────────────
\solutiontitle{3}{46}{\algo\ Algorithme~: simplifier une fraction}

Trace pour $p=36$, $q=48$~:
PGCD par Euclide~: $48=1\times36+12$~; $36=3\times12+0$~: $d=12$.
Fraction simplifi\'ee~: $36/12=3$ et $48/12=4$~: $\mathbf{3/4}$.

Algorithme Euclide int\'egr\'e~:
\begin{algorithm}[H]
\caption{PGCD + simplification}
\Lire $p, q$\;
$a\leftarrow p$~; $b\leftarrow q$\;
\TantQue{$b\ne0$}{$r\leftarrow a\bmod b$~; $a\leftarrow b$~; $b\leftarrow r$\;}
$d\leftarrow a$\;
\Afficher $p/d$, \og/\fg, $q/d$\;
\end{algorithm}

Python~: \texttt{from math import gcd; print(p//gcd(p,q), q//gcd(p,q))}.

%────────────────────────────────────────────────────────────────
\solutiontitle{3}{47}{\algo\ Algorithme~: type d'un r\'eel}

\begin{algorithm}[H]
\caption{Classifier un r\'eel}
\Lire $x$\;
\Si{$x = \lfloor x\rfloor$}{\Afficher \og entier\fg\;}
\SinonSi{les d\'ecimales se r\'ep\`etent (test approximatif)}{\Afficher \og rationnel non entier\fg\;}
\Sinon{\Afficher \og possiblement irrationnel\fg\;}
\end{algorithm}

Tests~: $x=0{,}5=1/2$~: rationnel.\quad
$x\approx1{,}41421\approx\sqrt{2}$~: possiblement irrationnel.\quad
$x=3$~: entier.

\emph{Remarque~:} un ordinateur ne peut pas d\'ecider avec certitude si un r\'eel est irrationnel
-- il ne voit qu'une approximation finie.

%────────────────────────────────────────────────────────────────
\solutiontitle{3}{48}{\algo\ Somme d'une s\'erie g\'eom\'etrique}

Trace pour $r=2$, $n=4$~:
\begin{center}
\renewcommand{\arraystretch}{1.2}
\begin{tabular}{cccc}
\toprule
$k$ & $t$ (avant) & $S$ (apr\`es $S\leftarrow S+t$) & $t$ (apr\`es $t\leftarrow t\times r$) \\
\midrule
0 & 1 & 1 & 2 \\
1 & 2 & 3 & 4 \\
2 & 4 & 7 & 8 \\
3 & 8 & 15 & 16 \\
4 & 16 & 31 & 32 \\
\bottomrule
\end{tabular}
\end{center}

$S=31=1+2+4+8+16$. V\'erification~: $\dfrac{2^5-1}{2-1}=31$.\checkmark

Version modifi\'ee avec arr\^et d\`es $S>1000$~: ici $S=1023$ apr\`es $k=9$.

%────────────────────────────────────────────────────────────────
\solutiontitle{3}{49}{\algo\ Dichotomie pour $\sqrt{a}$}

\begin{algorithm}[H]
\caption{$\sqrt{a}$ par dichotomie}
\Lire $a$~;\quad $g\leftarrow0$~;\quad $d\leftarrow\max(1,a)$~;\quad $\varepsilon\leftarrow10^{-4}$\;
\TantQue{$d-g>\varepsilon$}{
  $m\leftarrow(g+d)/2$\;
  \Si{$m^2\leq a$}{$g\leftarrow m$\;}
  \Sinon{$d\leftarrow m$\;}
}
\Afficher $(g+d)/2$\;
\end{algorithm}

Nombre d'it\'erations pour atteindre $\varepsilon=10^{-4}$~:
$\lceil\log_2(a/\varepsilon)\rceil$ -- environ $13$--$17$ it\'erations pour $a\leq100$.

Comparaison avec H\'eron~: H\'eron atteint $10^{-8}$ en $\approx5$ it\'erations
(convergence quadratique vs logarithmique pour la dichotomie).

\subsection*{Probl\`emes}
\addcontentsline{toc}{subsection}{Probl\`emes}

%────────────────────────────────────────────────────────────────
\psolutiontitle{3}{1}{Mesures et ordres de grandeur}

\begin{enumerate}
  \item $1\text{ ann\'ee}\approx3{,}15\times10^7\text{ s}$~;
        $1\text{ si\`ecle}=3{,}15\times10^9\text{ s}$.
        En notation scientifique~: $3{,}15\times10^7$ et $3{,}15\times10^9$.
  \item Vitesse de la lumi\`ere~: $c=3\times10^8\text{ m/s}$.
        Diam\`etre de la Voie Lact\'ee~: $d\approx9{,}5\times10^{20}\text{ m}$.
        Temps de travers\'ee~: $d/c\approx3{,}17\times10^{12}\text{ s}$
        $\approx10^5\text{ ans}$.
  \item Masse d'un \'electron~: $9{,}1\times10^{-31}\text{ kg}$.
        $10^{30}$ \'electrons~: $9{,}1\times10^{-31}\times10^{30}=0{,}91\text{ kg}$.
        Moins de $1\text{ kg}$.
\end{enumerate}

%────────────────────────────────────────────────────────────────
\psolutiontitle{3}{2}{Encadrements en physique}

\begin{enumerate}
  \item $g=9{,}81\text{ m/s}^2$. Lors d'une chute de $h=5\text{ m}$~:
        $v=\sqrt{2gh}=\sqrt{2\times9{,}81\times5}=\sqrt{98{,}1}\approx9{,}9\text{ m/s}$.
        Encadrement~: $9^2=81<98{,}1<100=10^2$, donc $9<v<10\text{ m/s}$.
  \item $T=2\pi\sqrt{L/g}$. Pour $L=1\text{ m}$~:
        $T=2\pi\sqrt{1/9{,}81}\approx2\pi\times0{,}319\approx2{,}006\text{ s}$.
        Encadrement~: $2<T<3$.
  \item Les valeurs mesur\'ees encadrent la vraie valeur~: si on mesure $v=10{,}2\pm0{,}5\text{ m/s}$,
        alors $9{,}7\leq v_\text{vraie}\leq10{,}7$.
\end{enumerate}

%────────────────────────────────────────────────────────────────
\psolutiontitle{3}{3}{Rationnels proches d'un irrationnel}

\begin{enumerate}
  \item Fractions approchant $\sqrt{2}$~: $1/1, 3/2, 7/5, 17/12, 41/29, 99/70$.
        Ces fractions viennent de l'algorithme des convergentes de la fraction
        continue $\sqrt{2}=[1;2,2,2,\ldots]$.
  \item $|p/q-\sqrt{2}|<1/q^2$ pour toute convergente (th\'eor\`eme de Hurwitz).
        Pour $p/q=99/70$~: $|99/70-\sqrt{2}|=|1{,}41429\ldots-1{,}41421\ldots|\approx8\times10^{-5}$
        et $1/70^2\approx2\times10^{-4}$~: $8\times10^{-5}<2\times10^{-4}$.\checkmark
  \item Ces approximations sont les \emph{meilleures rationnelles} pour le d\'enominateur donn\'e.
\end{enumerate}

%────────────────────────────────────────────────────────────────
\psolutiontitle{3}{4}{In\'egalit\'e de Bernoulli}

\textbf{Th\'eor\`eme} (Bernoulli)~: pour $x\geq-1$ et $n\in\N$~:
$(1+x)^n\geq1+nx$.

\begin{enumerate}
  \item $n=1$~: $(1+x)^1=1+x=1+1\cdot x$.\checkmark
  \item $n=2$~: $(1+x)^2=1+2x+x^2\geq1+2x$ car $x^2\geq0$.\checkmark
  \item $n=3$~: $(1+x)^3=1+3x+3x^2+x^3$.
        Pour $x\geq-1$~: $3x^2+x^3=x^2(3+x)\geq0$ (car $3+x\geq2>0$).
        Donc $(1+x)^3\geq1+3x$.\checkmark
  \item Application~: $(1{,}1)^{10}=(1+0{,}1)^{10}\geq1+10\times0{,}1=2$.
        En r\'ealit\'e $(1{,}1)^{10}\approx2{,}59>2$.\checkmark
\end{enumerate}

%────────────────────────────────────────────────────────────────
\psolutiontitle{3}{5}{La spirale de Pythagore}

La spirale de Pythagore construit $\sqrt{2},\sqrt{3},\sqrt{4},\ldots$ par triangles rectangles.

\begin{enumerate}
  \item Triangle $T_1$~: jambes $1$ et $1$, hypot\'enuse $\sqrt{2}$.
        Triangle $T_2$~: jambes $\sqrt{2}$ et $1$, hypot\'enuse $\sqrt{3}$.
        En g\'en\'eral~: $T_n$ a des jambes $\sqrt{n}$ et $1$, hypot\'enuse $\sqrt{n+1}$.
  \item V\'erification~: $(\sqrt{n})^2+1^2=n+1=(\sqrt{n+1})^2$.\checkmark\hfill$\square$
  \item $\sqrt{n}$ est rationnelle si et seulement si $n$ est un carr\'e parfait.
        La spirale produit notamment
        \[\sqrt{2},\sqrt{3},\sqrt{5},\sqrt{6},\sqrt{7},\sqrt{8},\sqrt{10},\ldots\]
        Ces nombres sont irrationnels lorsque $n$ n'est pas un carr\'e parfait.
  \item Le nombre de triangles pour atteindre $\sqrt{N}$ est $N-1$.
\end{enumerate}

%────────────────────────────────────────────────────────────────
\psolutiontitle{3}{6}{Valeur absolue et in\'equations physiques}

\begin{enumerate}
  \item $|T-20|\leq5\Leftrightarrow15\leq T\leq25$. Zone de confort~: $[15^{\circ}\text{C},25^{\circ}\text{C}]$.
  \item $|v-v_0|\leq\delta$~: la vitesse est dans $[v_0-\delta,v_0+\delta]$.
        Pour $v_0=50\text{ km/h}$, $\delta=5$~: $[45,55]$.
  \item $|f-440|\leq10$~: fr\'equence dans $[430,450]$ Hz (accord du La).
  \item En g\'en\'eral, $|x-a|\leq r$ repr\'esente un intervalle de \emph{tol\'erance}
        centr\'e en $a$ de \emph{rayon} $r$.
\end{enumerate}

%────────────────────────────────────────────────────────────────
\psolutiontitle{3}{7}{Encadrements et d\'ecimales de $\sqrt{2}$}

\begin{enumerate}
  \item $1{,}4^2=1{,}96<2<2{,}25=1{,}5^2$~: $1{,}4<\sqrt{2}<1{,}5$.\hfill$\square$
  \item $1{,}41^2=1{,}9881<2<1{,}9984=1{,}412^2$... $1{,}42^2=2{,}0164$~:
        $1{,}41<\sqrt{2}<1{,}42$.\hfill$\square$
  \item $1{,}414^2=1{,}999396<2<1{,}4149...^2$~:
        $1{,}414<\sqrt{2}<1{,}415$.\hfill$\square$
  \item Par dichotomie ou calcul~: $\sqrt{2}\approx1{,}41421356\ldots$
  \item Chaque fois qu'on affine le pas d'un facteur~$10$, on gagne une d\'ecimale.
        C'est la convergence lin\'eaire ($O(1/10^n)$).
\end{enumerate}

%────────────────────────────────────────────────────────────────
\psolutiontitle{3}{8}{Organisation et classement des fractions}

\begin{enumerate}
  \item Pour comparer $p/q$ et $r/s$ ($q,s>0$)~: $p/q<r/s\Leftrightarrow ps<rq$.
  \item Classer $1/3$, $2/5$, $3/8$, $5/12$~:
        $1/3\approx0{,}333$~; $2/5=0{,}400$~; $3/8=0{,}375$~; $5/12\approx0{,}417$.
        Ordre~: $1/3<3/8<2/5<5/12$.
  \item M\'ediane de Farey~: entre $p/q$ et $r/s$, la fraction $(p+r)/(q+s)$ (m\'ediante)
        est entre les deux et est la fraction de plus petit d\'enominateur dans cet intervalle.
        Exemple~: m\'ediante de $1/3$ et $2/5$ est $3/8$.\checkmark
\end{enumerate}

%────────────────────────────────────────────────────────────────
\psolutiontitle{3}{9}{Distance et m\'ediane}

\begin{enumerate}
  \item $d(A,B)=|b-a|$. Pour minimiser $\sum|x_i-c|$ (somme des distances),
        $c$ est la \textbf{m\'ediane} de la s\'erie $\{x_i\}$.
  \item Preuve pour $n=3$ valeurs $x_1\leq x_2\leq x_3$~:
        $f(c)=|c-x_1|+|c-x_2|+|c-x_3|$ est minimis\'ee en $c=x_2$ (m\'ediane).
        Pour $c=x_2$~: $f(x_2)=(x_2-x_1)+(x_3-x_2)=x_3-x_1$.
        Pour tout autre $c$, $f(c)\geq x_3-x_1$.\hfill$\square$
  \item Contrairement \`a la moyenne, la m\'ediane minimise la somme des
        \emph{distances} (et non des carr\'es des distances).
\end{enumerate}

%────────────────────────────────────────────────────────────────
\psolutiontitle{3}{10}{Notation scientifique et astronomie}

\begin{enumerate}
  \item Distance Terre-Lune~: $3{,}84\times10^8\text{ m}$.
        \`A $c=3\times10^8\text{ m/s}$~: $t=3{,}84\times10^8/(3\times10^8)\approx1{,}28\text{ s}$.
  \item Distance Terre-Soleil~: $1{,}50\times10^{11}\text{ m}$.
        $t=1{,}50\times10^{11}/(3\times10^8)\approx500\text{ s}\approx8{,}3\text{ min}$.
  \item \'Etoile la plus proche (Proxima Centauri)~: $4{,}24$ ann\'ees-lumi\`ere
        $=4{,}24\times3{,}15\times10^7\times3\times10^8\approx4\times10^{16}\text{ m}$.
        $t\approx4{,}24\text{ ans}$.
\end{enumerate}

%────────────────────────────────────────────────────────────────
\psolutiontitle{3}{11}{Fractions et th\`eme harmonique}

\begin{enumerate}
  \item Harmonie musicale~: octave $=2/1$, quinte $=3/2$, quarte $=4/3$.
        Les intervalles sont des rapports de fractions simples.
  \item S\'erie harmonique~: $H_n=1+1/2+1/3+\cdots+1/n$.
        $H_1=1$, $H_2=3/2$, $H_4=2{,}083$, $H_{10}\approx2{,}93$.
        $H_n\to+\infty$ mais tr\`es lentement.
  \item Moyenne harmonique de $a$ et $b$~: $H=\dfrac{2ab}{a+b}$.
        Pour $a=3$, $b=6$~: $H=\dfrac{36}{9}=4$.
        Ordre~: $H\leq G\leq A$ (harmonique $\leq$ g\'eom\'etrique $\leq$ arithm\'etique).
\end{enumerate}

%────────────────────────────────────────────────────────────────
\psolutiontitle{3}{12}{D\'ecrire des ensembles de r\'eels}

\begin{enumerate}
  \item $\{x\in\R\mid x^2\leq4\}=[-2,2]$ (car $|x|\leq2$).
  \item $\{x\in\R\mid|x-3|<2\}=]1,5[$.
  \item $\{x\in\R\mid x^2>x\}=\{x\mid x(x-1)>0\}=]-\infty,0[\cup]1,+\infty[$.
  \item $\{x\in\R\mid\sqrt{x-1}\leq2\}=\{x\mid x-1\in[0,4]\}=[1,5]$.
  \item $\{x\in\R^*\mid1/x\in]0,1]\}=\{x>0\mid x\geq1\}=[1,+\infty[$.
\end{enumerate}

%────────────────────────────────────────────────────────────────
\psolutiontitle{3}{13}{\logiq\ Logique des ensembles de nombres}

\begin{enumerate}
  \item Contre-exemple~: $x=\sqrt{2}$ est irrationnel, mais $x^2=2\in\Q$.
        Donc $x^2\in\Q\not\Rightarrow x\in\Q$.
  \item Si $x=p/q$ et $y=r/s$ ($q,s\ne0$)~:
        $x+y=\dfrac{ps+rq}{qs}\in\Q$.\hfill$\square$
  \item Supposons $x\in\Q$ et $y\notin\Q$.
        Si $x+y\in\Q$~: $y=(x+y)-x\in\Q$ (diff\'erence de deux rationnels)~:
        contradiction. Donc $x+y\notin\Q$.\hfill$\square$
  \item Fausse. Contre-exemple~: $x=\sqrt{2}$, $y=\sqrt{2}$~:
        $xy=2\in\Q$ mais $x,y\notin\Q$.
\end{enumerate}

%────────────────────────────────────────────────────────────────
\psolutiontitle{3}{14}{\logiq\ Logique et valeur absolue}

\begin{enumerate}
  \item $\forall x\in\R, |x|\geq0$. Vrai par d\'efinition.\hfill$\square$
  \item $|x+y|\leq|x|+|y|$ (in\'egalit\'e triangulaire)~: voir exo~26.\hfill$\square$
  \item R\'eciproque de l'in\'egalit\'e triangulaire~:
        $|x|+|y|=|x+y|\Leftrightarrow x$ et $y$ de m\^eme signe (ou l'un nul).
        La r\'eciproque \og si $|x|+|y|=|x+y|$ alors $xy\geq0$\fg\ est vraie.\hfill$\square$
\end{enumerate}

%────────────────────────────────────────────────────────────────
\psolutiontitle{3}{15}{\algo\ Trier et calculer la m\'ediane}

\begin{algorithm}[H]
\caption{Tri insertion + m\'ediane}
\KwIn{liste $L$ de $n$ valeurs}
Trier $L$ par insertion (voir Ch.~2)\;
\Si{$n \bmod 2 = 1$}{\Retourner $L[\lfloor n/2\rfloor+1]$\;}
\Sinon{\Retourner $(L[n/2]+L[n/2+1])/2$\;}
\end{algorithm}

Test sur $\{5,2,8,1,9,3\}$ ($n=6$)~:
Tri~: $[1,2,3,5,8,9]$.
M\'ediane~: $(3+5)/2=4$.

En Python~:
\begin{lstlisting}[style=python]
L = sorted([5, 2, 8, 1, 9, 3])
n = len(L)
med = L[n//2] if n%2==1 else (L[n//2-1]+L[n//2])/2
print(f"Mediane : {med}")
\end{lstlisting}

%────────────────────────────────────────────────────────────────
\psolutiontitle{3}{16}{\algo\ D\'ecimales d'un irrationnel}

\begin{algorithm}[H]
\caption{Encadrement de $\sqrt{2}$ \`a $10^{-k}$}
$a\leftarrow1$~;\quad $b\leftarrow2$\;
\Pour{$k\leftarrow1$ \KwTo $5$}{
  $h\leftarrow10^{-k}$~;\quad $x\leftarrow a$\;
  \TantQue{$x^2\leq2$}{$x\leftarrow x+h$\;}
  $b\leftarrow x$~;\quad $a\leftarrow x-h$\;
  \Afficher $k$, $a$, $b$, \og pr\'ecision~: \fg, $h$\;
}
\end{algorithm}

R\'esultats~:
$k=1$~: $1{,}4<\sqrt{2}<1{,}5$~;
$k=2$~: $1{,}41<\sqrt{2}<1{,}42$~;
$k=3$~: $1{,}414<\sqrt{2}<1{,}415$~;
$k=4$~: $1{,}4142<\sqrt{2}<1{,}4143$~;
$k=5$~: $1{,}41421<\sqrt{2}<1{,}41422$.

\subsection*{D\'efis}
\addcontentsline{toc}{subsection}{D\'efis}

\noindent\phantomsection
{\color{BO}\bfseries\small D\'efi~1~\textendash\ \textit{Approximation de $\sqrt{2}$ par l'algorithme babylonien}}
\par\smallskip

\begin{enumerate}
  \item L'algorithme babylonien pose $u_{n+1}=(u_n+2/u_n)/2$ (m\'ethode de H\'eron).
  \item Trace pour $u_0=2$~:
        $u_1=(2+1)/2=1{,}5$~;
        $u_2=(1{,}5+4/3)/2\approx1{,}4167$~;
        $u_3\approx1{,}4142156$.
        Apr\`es $4$ it\'erations~: $8$ d\'ecimales correctes.
  \item Convergence \textbf{quadratique}~: le nombre de d\'ecimales correctes double \`a
        chaque it\'eration.
  \item D\'emonstration~: $u_n-\sqrt{2}=\dfrac{(u_{n-1}-\sqrt{2})^2}{2u_{n-1}}$.
        Si $e_n=u_n-\sqrt{2}$~: $e_n\approx e_{n-1}^2/(2\sqrt{2})$.\hfill$\square$
\end{enumerate}

\noindent\phantomsection
{\color{BO}\bfseries\small D\'efi~2~\textendash\ \textit{Le nombre d'or $\varphi$}}
\par\smallskip

\begin{enumerate}
  \item $\varphi=\dfrac{1+\sqrt{5}}{2}$ est solution de $x^2-x-1=0$.\hfill$\square$
  \item $\varphi^2=\varphi+1$~: $\left(\dfrac{1+\sqrt{5}}{2}\right)^2=\dfrac{6+2\sqrt{5}}{4}=\dfrac{3+\sqrt{5}}{2}=\dfrac{1+\sqrt{5}}{2}+1=\varphi+1$.\checkmark
  \item $1/\varphi=\varphi-1\approx0{,}618$. En effet~: $1/\varphi=\dfrac{2}{1+\sqrt{5}}=\dfrac{2(1-\sqrt{5})}{(1-5)}=\dfrac{\sqrt{5}-1}{2}=\varphi-1$.
  \item Fibonacci~: $F_{n+1}/F_n\to\varphi$ (voir Ch.~2, exercice~30).
  \item Fraction continue~: $\varphi=1+\cfrac{1}{1+\cfrac{1}{1+\cfrac{1}{\ddots}}}=[1;1,1,1,\ldots]$.
        C'est la fraction continue la plus lentement convergente (nombre le plus irrationnel).
\end{enumerate}

\noindent\phantomsection
{\color{BO}\bfseries\small D\'efi~3~\textendash\ \textit{Sommes et produits de rationnels et d'irrationnels}}
\par\smallskip

\begin{center}
\renewcommand{\arraystretch}{1.4}
\begin{tabular}{|l|c|c|}
\hline
\rowcolor{BN!10}
Op\'eration & R\'esultat & Exemple \\
\hline
$\Q+\Q$ & $\in\Q$ & $1/2+1/3=5/6$ \\
$\Q\times\Q$ & $\in\Q$ & $2/3\times3/4=1/2$ \\
$\R\setminus\Q+\R\setminus\Q$ & Peut \^etre $\Q$ ou non & $\sqrt{2}+(-\sqrt{2})=0$ \\
$\R\setminus\Q\times\R\setminus\Q$ & Peut \^etre $\Q$ ou non & $\sqrt{2}\times\sqrt{2}=2$ \\
$\Q+(\R\setminus\Q)$ & $\in\R\setminus\Q$ & $1+\sqrt{2}\notin\Q$ \\
$\Q^*\times(\R\setminus\Q)$ & $\in\R\setminus\Q$ & $2\sqrt{2}\notin\Q$ \\
\hline
\end{tabular}
\end{center}

D\'emonstration de la derni\`ere ligne~: si $r\in\Q^*$ et $x\notin\Q$,
supposons $rx=s\in\Q$. Alors $x=s/r\in\Q$ (quotient de rationnels)~:
contradiction.\hfill$\square$

\noindent\phantomsection
{\color{BO}\bfseries\small D\'efi~4~\textendash\ \textit{Repr\'esentation d\'ecimale des rationnels}}
\par\smallskip

\begin{enumerate}
  \item Tout rationnel $p/q$ a une repr\'esentation d\'ecimale p\'eriodique.
        La p\'eriode est de longueur $\leq q-1$ (valeur du reste dans l'algorithme de division).
  \item $1/7=0{,}\overline{142857}$ (p\'eriode $6$~; $1/7$ a $q=7$).
  \item $1/6=0{,}1\overline{6}$ (p\'eriode~1 apr\`es le d\'ecalage).
  \item R\'eciproque~: tout d\'ecimal p\'eriodique est rationnel.
        Preuve~: si $x=0{,}\overline{a_1\ldots a_k}$~, alors $10^k x-x=(10^k-1)x=a_1\ldots a_k\in\Z$,
        d'o\`u $x=n/(10^k-1)\in\Q$.\hfill$\square$
\end{enumerate}

\noindent\phantomsection
{\color{BO}\bfseries\small D\'efi~5~\textendash\ \textit{Construction de $\sqrt{2}$ \`a la r\`egle et au compas}}
\par\smallskip

\begin{enumerate}
  \item Construire un triangle rectangle isoc\`ele de jambes $1$~:
        l'hypot\'enuse mesure $\sqrt{2}$ (th\'eor\`eme de Pythagore).
  \item Construire par induction~: $\sqrt{n+1}$ depuis $\sqrt{n}$~:
        triangle rectangle de jambes $\sqrt{n}$ et $1$,
        hypot\'enuse $\sqrt{n+1}$ (spirale de Pythagore).
  \item Pour $\sqrt{p/q}$~: construire $\sqrt{p}$ et $\sqrt{q}$,
        puis utiliser la mean proportionnelle.
  \item Les constructibles \`a la r\`egle et au compas sont exactement
        les r\'eels obtenus par $+,-,\times,\div,\sqrt{\cdot}$ \`a partir
        de $\Q$.
\end{enumerate}

\noindent\phantomsection
{\color{BO}\bfseries\small D\'efi~6~\textendash\ \textit{Calcul alg\'ebrique avanc\'e avec les radicaux}}
\par\smallskip

\begin{enumerate}
  \item $(\sqrt{2}+\sqrt{3})^4=(5+2\sqrt{6})^2=25+20\sqrt{6}+24=49+20\sqrt{6}$.
  \item $\dfrac{1}{\sqrt{2}+\sqrt{3}+\sqrt{5}}=\dfrac{\sqrt{2}+\sqrt{3}-\sqrt{5}}{(\sqrt{2}+\sqrt{3})^2-5}=\dfrac{\sqrt{2}+\sqrt{3}-\sqrt{5}}{2\sqrt{6}}=\dfrac{(\sqrt{2}+\sqrt{3}-\sqrt{5})\sqrt{6}}{12}$.
  \item $\sqrt[3]{2+\sqrt{5}}+\sqrt[3]{2-\sqrt{5}}$~: on v\'erifie que si $t$ est cette expression,
        alors $t^3=4-3t$ (en \'ecrivant $t=a+b$ avec $ab=\sqrt[3]{4-5}=\sqrt[3]{-1}=-1$),
        donc $t^3+3t-4=0$~: solution $t=1$.
  \item $(\sqrt{3}+\sqrt[3]{2})^2\notin\Q$~: l'\'el\'evation au carr\'e ne simplifie pas.
        $3+2^{2/3}+2\sqrt{3}\cdot2^{1/3}$~: contient des termes de degr\'es diff\'erents.
\end{enumerate}

\noindent\phantomsection
{\color{BO}\bfseries\small D\'efi~7~\textendash\ \textit{\logiq\ Logique et irrationalit\'e de $\sqrt{p}$}}
\par\smallskip

\begin{enumerate}
  \item \textbf{Lemme.} Si $p$ est premier et $p\mid n^2$, alors $p\mid n$.
        Preuve~: si $p\nmid n$, alors $\pgcd(p,n)=1$ (car $p$ est premier),
        et par le lemme de Gauss $p\nmid n^2$~: contrapos\'ee.\hfill$\square$
  \item \textbf{Th\'eor\`eme.} $\sqrt{p}\notin\Q$ pour tout premier $p$.
        Supposons $\sqrt{p}=a/b$ irr\'eductible.
        $p b^2=a^2 \Rightarrow p\mid a \Rightarrow a=pc \Rightarrow p^2c^2=pb^2
        \Rightarrow p\mid b$~: contradiction.\hfill$\square$
  \item G\'en\'eralisation~: $\sqrt{n}\notin\Q$ si et seulement si
        $n$ n'est \textbf{pas} un carr\'e parfait.
        \'Enonc\'e logique~: $\sqrt{n}\in\Q \Leftrightarrow \exists k\in\N, n=k^2$.
\end{enumerate}

\noindent\phantomsection
{\color{BO}\bfseries\small D\'efi~8~\textendash\ \textit{\algo\ Algorithme de Newton-H\'eron g\'en\'eralis\'e}}
\par\smallskip

\begin{enumerate}
  \item Pour $k=2$~: $u_{n+1}=\dfrac{u_n+a/u_n}{2}$~: c'est la m\'ethode de H\'eron.\checkmark
  \item Algorithme g\'en\'eral~: $u_{n+1}=\dfrac{(k-1)u_n+a/u_n^{k-1}}{k}$.
        Arr\^eter quand $|u_{n+1}-u_n|<\varepsilon$.
  \item Pour $k=2$, $a=2$~: convergence en $5$ it\'erations \`a $10^{-8}$.
        Pour $k=3$, $a=2$~: $\sqrt[3]{2}$~; convergence similaire.
  \item Preuve de $u_n>\sqrt[k]{a}$ si $u_0>\sqrt[k]{a}$~:
        par l'AM--GM, $\dfrac{(k-1)u+a/u^{k-1}}{k}\geq \sqrt[k]{u^{k-1}\cdot(a/u^{k-1})}=a^{1/k}$.
        D'o\`u $u_{n+1}\geq a^{1/k}$, avec \'egalit\'e ssi $u_n=a^{1/k}$.\hfill$\square$
\end{enumerate}


\needspace{3\baselineskip}
\needspace{3\baselineskip}
